\chapter{Chain 细节全拆：量化、欧拉环、并行抢第一刀}
\label{ch:chain-details}

\section{量化 $Q(b)$：把字节磨粗糙}

$Q(b)=b\gg q$。$q=0$ 精细；$q$ 变大则许多字节挤进同一颜色桶。

\begin{figure}[htbp]
\centering
\begin{tikzpicture}[font=\small]
  \node[bytecell] (b) {0x97};
  \node[softbox=Gen4Color, right=1.0cm of b] (q0) {$q=0$\\键 0x97};
  \node[softbox=Gen4Color, right=0.6cm of q0] (q2) {$q=2$\\键 0x25};
  \draw[-{Stealth}, thick] (b)--(q0);
  \draw[-{Stealth}, thick] (b)--(q2);
\end{tikzpicture}
\caption{quant\_q 合法 0..7；存超块。}
\label{fig:quant}
\end{figure}

\section{LFSR 拧锁逐步}

\begin{toybox}[4 位示意]
$S=0101$，来键 3：左旋 $\to$ 1010，异或 0011 $\to$ 1001。\\
若 $L=2$（低 2 位），要对齐需 $**00$。继续拧直到对齐。
\end{toybox}

真名：\texttt{UpdateChainSignature} $=\mathrm{rotl}_1\oplus Q$；\texttt{ChainStrideMask}$=2^L-1$，$L\in[8,22]$。

\section{Limit Push：对齐之间不看珠}

\begin{metaphor}[高速公路服务区]
非对齐：油门 LFSR（每字节 $O(1)$）。\\
对齐：才进服务区检查骨架（$O(w)$）和环。\\
若每字节都检查，Chain 就失去相对 Gear 的优势。
\end{metaphor}

\section{$\hat\beta_1$：边太多出现环}

\begin{equation}
\hat\beta_1=\max(0,\,E-V+C)
\end{equation}
$V=$出现过的颜色数，$E=$去重无向异色边，$C=$连通块数。

\begin{figure}[htbp]
\centering
\begin{tikzpicture}[font=\small]
  \node[align=left, softbox=Gen4Color, minimum width=5.5cm]
    {线状：边少，$E-V+C\le 0$\\环状：边多，$E-V+C>0$ 过闸};
\end{tikzpicture}
\caption{参考实现 vs bitset 快路径：答案必须 bit-identical。}
\label{fig:euler}
\end{figure}

\section{beta1 开关与 MVP 升级禁令}

新仓默认开环检。老 MVP 仓没有这 bit —— \textbf{禁止就地拧开}，必须新开房重切。\\
形象：旧门锁芯不同，不能只换钥匙齿。

\section{并行 256KB：预热锁芯，抢最左刀}

\begin{figure}[htbp]
\centering
\begin{tikzpicture}[font=\footnotesize]
  \foreach \i in {0,1,2,3} {
    \node[softbox=Gen4Color, minimum width=2.3cm] (s\i) at (\i*2.7,0) {段\i};
  }
  \node[font=\small, above] at (4,1.0) {先串行预热各段起点的 $S$};
  \draw[-{Stealth}, thick, gray] (s0)--(s1)--(s2)--(s3);
  \node[softbox=CutRed, below=1.1cm of s1] (cut) {段1 先找到};
  \draw[-{Stealth}, CutRed, thick] (s1)--(cut);
  \node[font=\footnotesize, gray, below=0.35cm of cut] {喊停右侧；不认更右的刀};
\end{tikzpicture}
\caption{\texttt{PARALLEL\_SCAN=1}；有跨 feed carry 时强制串行。}
\label{fig:chain-par}
\end{figure}

\section{Resume 带走什么}

Hom 可不带珠串；Chain 常带走锁状态 $S$ 和 keys 窗 —— 不然跨盒子要重拧很长一段。

\section{标定 \texttt{CalibrateChainStrideLog}}

在 1MB 样本上试 $L$（含 beta1），使平均块长落入 $[0.85,1.15]\times\mathrm{avg}$。\\
$L$ 写入超块 \texttt{reserved[0:1]}。

\section{细节对照表}

\begin{center}
\small
\begin{tabular}{@{}lp{7.2cm}@{}}
\toprule
技术点 & 形象说法 \\
\midrule
\texttt{ChainQuantKey} & 磨粗糙 \\
\texttt{ChainPrimaryDelta} & 与 Hom 同款三边 \\
\texttt{ChainBeta1Gf2} & 有没有打结 \\
\texttt{InitChainSignature} & 开窗灌满锁芯 \\
\texttt{force\_serial} & 过盒子边界别抢跑 \\
切点冻结 & 优化只能换发动机外壳，不能改交规 \\
\bottomrule
\end{tabular}
\end{center}
